\chapter{Reduction of Endomorphisms}

\section{Principal Ideal Domain}

\begin{definitionenv}
    Let $R$ be a commutative unitary ring.
    By definition, an ideal of $R$ is a sub-$R$-module of $R$. 
    If an ideal $I$ of $R$ is generated by one element, we say that $I$ is a \textbf{principal ideal}.

    \quad If all ideals of $R$ are principal, we say that $R$ is a \textbf{principal ideal ring}.

    \quad Let $a,b \in R$, then $aR \subseteq bR$ if and only if there exists $c \in R$ such that $a = b c$.
    We denote this condition as $b \mid a$, and say that $b$ is a \textbf{divisor} of $a$ and $a$ is a \textbf{multiple} of $b$.
\end{definitionenv}

\begin{definitionenv}\label{prop:6.1.3}
    Let $R$ be a commutative unitary ring.
    If the neutral element for the addition and multiplication are different, and if
    $$\forall s,t \in R\setminus\{0\},\ st \in R\setminus \{0\},$$
    then we say that $R$ is an \textbf{integral domain}.
\end{definitionenv}

\begin{propositionenv}
    Let $R$ be an integral domain and $a,b \in R$.
    Then $aR = bR$ if and only if there exists $c \in R^\times$ such that $a = b c$.
\end{propositionenv}
\begin{proofenv}
    We assume that there exists $c \in R^\times$ such that $a = b c$.
    By definition, $aR \subseteq bR$.
    Since $c$ is invertible, also $bR \subseteq aR$.

    \quad Conversely, assume that $aR = bR$, then we have $x,y \in R$ such that $a = b x$, $b = a y$.
    If $a = 0$, then $b = 0$.
    In this case, $a = 1 \cdot b$.
    If $a \neq 0$, then $a = b x = a y x$.
    So $y x = 1$ and $x y \in R^\times$.
\end{proofenv}

\begin{definitionenv}
    Let $K$ be a commutative unitary ring.
    Recall that a non-zero polynomial $P \in K[T]$ can be written as
    $$a_0 + a_1 T +\cdots + a_n T^n,$$
    where $a_0,\cdots a_n \in K$ and $a_n \neq 0$.
    
    \quad The natural number $n$ is called the \textbf{degree} of $P$, and denote as $\deg (P) \in \NN$.
    The element $a_n$ is called the \textbf{leading coefficient} of $P$, and denote as $c(P) \in K$.

    \quad If $c(P) = 1$, then we say that $P$ is a \textbf{monic}.

    \quad By convention, the degree of $0$ is defined to be $-\infty$,
    and the zero polynomial is assumed to be monic.
\end{definitionenv}

\begin{propositionenv}
    Let $K$ be an integral domain.
    Then the ring $K[T]$ is also a integral domain.
    It is a principal ideal ring if $K$ is a field.
\end{propositionenv}
\begin{proofenv}
    Since $K$ is an integral domain, if $P,Q \in K[T]$ are non-zero in $K[T]$, then
    $$\deg (PQ) = \deg (P) + \deg (Q),\quad c(PQ) = c(P)c(Q).$$
    Hence $K[T]$ is also an integral domain.

    \quad We check that $K[T]$ is a principal ideal ring when $K$ is a field.
    We solve it using Euclidean division of polynomials.
    If $I$ is a zero ideal, then $(0) = I$.
    If $I$  is non-zero, then there exists non-zero $P \in I$ such that $\deg(P)$ is minimal inside $I$.
    For any $F \in K[T]$, there exists a unique pair $(Q,R) \in K[T]^2$ such that 
    $$F = PQ + R, \quad \deg(R) < \deg(P).$$
    Hence, since the degree of $P$ is minimal, $R = 0$, so $F = PQ$.
    So $I = (P)$.
\end{proofenv}

\begin{remark}
    Let $K$ be a field,
    and $I$ be an ideal of $K[T]$.
    Then there exists a unique monic $P \in I$ such that $I = (P)$.
\end{remark}

\begin{definitionenv}
    Let $A$ be an integral domain.
    We say that $a \in A\setminus A^\times$ is \textbf{irreducible} if $a$ can not be decomposed as a product of two elements in $A\setminus A^\times$.

    \quad In case when $A = K[T]$,
    an irreducible element is called an \textbf{irreducible polynomial}.
\end{definitionenv}

\begin{definitionenv}
    Let $A$ be a commutative unitary ring and $I$ be an ideal of $A$.
    We say that $I$ is \textbf{prime ideal} if the quotient ring $A/I$ is an integral domain.
    This is equivalent to the condition that $I \neq A$ and
    $$\forall a,b \in A \setminus I, ab \in A \setminus I.$$
    We say that $I$ is \textbf{maximal} if $A/I$ is a field.
\end{definitionenv}

\begin{propositionenv}\label{prop:6.1.9}
    Let $A$ be a commutative unitary ring.
    The maximal ideals of $A$ are precisely the maximal elements in the set of al ideals of $A$ ordered by inclusion $\subseteq$.
\end{propositionenv}
\begin{proofenv}
    Let $I$ be an ideal of $A$.
    Suppose that $A/I$ is a field.
    If $J$ is another ideal of $A$ such that $J \supsetneq I$,
    then there exists $a \in J \setminus I$ such that the class of $a$ in $A/I$ is non-zero.
    This implies that there exists $b \in A$ such that $ab - 1 \in I$.
    In particular,
    $$1 = (1 - ab) + ab \in J.$$
    Hence, $J = A$.

    \quad Conversely, suppose that $I \neq A$ is maximal amongst ideals of $I$ with respect to inclusion.
    If $a \in A \setminus I$ is such that the class of $a$ is not invertible in $A/I$, then the set
    $$J = \{ab + c \mid b \in A, c \in I\}$$
    is an ideal of $A$ containing $I$.
    $J \neq A$, since $1$ does not belongs to $J$.
    Contradiction with the maximality of $I$.
\end{proofenv}

\begin{propositionenv}
    Let $A$ be an integral domain, $a \in A$.
    If $(a)$ is a prime ideal,
    then $a$ is an irreducible element of $A$.
\end{propositionenv}
\begin{proofenv}
    Suppose that $a = bc$ with $b,c \in A\setminus A^\times$.
    Since $I = (a)$ is prime,
    then $b$ and $c$ can't be simultaneously belong to $A\setminus I$. 
    Suppose that $b \in I$,
    then there exists $d \in A$ such that
    $b = ad = bcd$.
    Hence $cd = 1$.
    This contradicts $c \in A \setminus A^\times$.
\end{proofenv}

\begin{propositionenv}
    Let $A$ be a principal ideal domain.
    If $a \in A$ is irreducible,
    then the ideal generated by $a$ is a maximal ideal.
\end{propositionenv}
\begin{proofenv}
    Let $I = (a)$.
    By proposition \ref{prop:6.1.9},
    it suffices to show that $I$ is not contained in any other ideal.

    \quad Suppose by contradiction that $I \subsetneq J \subsetneq A$.
    Since $A$ is a principal ideal domain,
    there exists $b \in A$ such that $J = (b)$.
    Since $I \subseteq J$,
    there exists $c \in A$ such that $a = bc$.
    Notice that $b \in A \setminus A^\times$, $c \in A \setminus A^\times$ since otherwise we would have $I = J$.
    Contradiction.
\end{proofenv}

\begin{definitionenv}
    Let $R$ be a commutative unitary ring.
    We say that $R$ is \textbf{noetherian} if for any increasing sequence
    $$I_0 \subseteq I_1 \subseteq I_2 \subseteq \cdots$$
    of ideals in $R$, there exists $n \in \NN$ such that
    $$I_n = I_{n+1} = I_{n+2}.$$
\end{definitionenv}

\begin{propositionenv}
    A commutative unitary ring $R$ is noetherian if and only if any ideal of $R$ is finitely generated.
    Any principal ideal domain is noetherian.
\end{propositionenv}
\begin{proofenv}
    Assume $R$ is noetherian.
    Let $I$ be an ideal of $R$.
    Assume that $I$ is not finitely generated.
    One can recursively construct a strictly increasing sequence
    $$I_0 \subseteq I_1 \subseteq \cdots I_n \subseteq I_{n+1} \subseteq \cdots$$
    using the infinite set of generations of $I$.
    Contradiction.
    Conversely assume that any $I$ in $R$ is finitely generated.
    Let 
    $$I_0 \subseteq I_1 \subseteq \cdots I_n \subseteq I_{n+1} \subseteq \cdots$$
    be an increasing sequence of ideals.
    Then $\displaystyle I \coloneq \bigcup_{n \in \NN} I_n$ is an ideal in $R$, so $I = (x_1,\cdots,x_p)$.
    Let $n \in \NN$ such that
    $$I_n \supseteq \{x_1,\cdots x_p\},$$
    hence $(x_1,\cdots,x_p) = I_n = I_{n+1} = \cdots$.
\end{proofenv}

\begin{corollaryenv}\label{cor:6.1.14}
    Let $R$ be a commutative unitary noetherian ring.
    Then any ideal $I \neq R$ is contained in some maximal ideal of $R$.
\end{corollaryenv}
\begin{proofenv}
    Suppose that $I$ is not contained in any maximal ideal, then recursively we obtain a strictly increasing sequence
    $$I \subseteq I_0 \subsetneq I_1 \subsetneq \cdots \subsetneq I_n \subsetneq I_{n+1} \subsetneq \cdots$$
    such that $I_n \neq R$ for any $n \in \NN$.
\end{proofenv}

\begin{remark}\label{6.1.15}
    The statement of corollary \ref{cor:6.1.14} is true for any commutative unitary ring under the axiom of choice.
    Let $R$ be an commutative unitary ring, $I \subsetneq R$ be an ideal.
    $$\Theta = \{\text{ideals of $J$ of } R \mid I \subseteq J \neq R\},$$
    $\Theta \neq \varnothing$ and ordered by $\subseteq$.
    If $\Theta_0 \neq \varnothing$, totally ordered subset of $\Theta$ then $\bigcup_{J \Theta_0} J$ is an upper bound of $\Theta_0$ in $\Theta$.
    By Zorn lemma, $\Theta$ admits a maximal element which is necessary a maximal ideal.
\end{remark}

\begin{propositionenv}\label{prop:6.1.16}
    Let $A$ be an principal ideal domain,
    then any $a \in A\setminus A^\times$ can be written as 
    $$a = p_1 \cdots p_n,$$
    where $p_1,\cdots, p_n \in A$ are irreducible.
    Moreover, if
    $$a = p_1 \cdots p_n = q_1 \cdots q_m,$$
    then $n = m$ and there exists $\sigma \in \mathfrak{S}_n$ such that $\forall i \in \{1,\cdots,n\}$,
    $p_i A = q_{\sigma(i)} A.$
\end{propositionenv}
\begin{proofenv}
    \begin{enumerate}
        \item \textit{Existence}
            We only need to consider the case where $a$ is not irreducible, by corollary \ref{cor:6.1.14},
            $aA = (a)$ is strictly contained in a maximal ideal of $A$.
            There exists $a_1 \in A \setminus A^\times$ such that
            $$a = p_1 a_1.$$
            Iterating this strategy, we obtain
            $$a = p_1 \cdots p_i a_i,$$
            where $p_1,\cdots, p_i$ are irreducible in $A$, $a_i \in A \setminus A^\times$.
            This procedure has to terminate, as otherwise we would construct a strictly increasing
            $$aA \subsetneq a_1 A \subsetneq \cdots \subsetneq a_n A \subsetneq \cdots.$$
            Hence, $a = p_1 \cdots p_n$.
        \item \textit{Uniqueness}
            We reason by induction on the minimal number of irreducible factors in the decomposition.
            If $a$ is irreducible, trivial.

            \quad If $a$ is not irreducible,
            assume now that $a = p_1 \cdots p_n$, $n \ge 2$ and $a$ can not be written as a product of $i < n$ irreducible elements.
            By induction hypothesis,
            we assume the uniqueness (up to permutation) for elements that can be written as a product of $i < n$ irreducible elements.
            If $a = q_1 \cdots q_m \in p_n A$,
            there exists $i\{1,\cdots,m\}$ such that $q_i \in p_n A$ (by our definition, we immediately know that a maximal ideal is prime).
            By permutation, without losing of generality, $i = m$, so $q_m A \subseteq p_n A$.
            Since $q_m$ is irreducible, $q_m A$ is maximal too.
            So $q_m A = p_n A$.
            By proposition \ref{prop:6.1.3},
            there exists $\lambda \in A^\times$ such that $q_m = \lambda p_n$.
            So $p_1\cdots p_{n-1} = (\lambda q_1) \cdots q_{m-1}.$
            Note that $q_1 A = (\lambda q_1) A$,
            we obtain the desired result using the induction hypothesis.
    \end{enumerate}
\end{proofenv}

\begin{remark}
    Let $K$ be a field and $M$ be the set of all monic irreducible polynomials in $K[T]$.
    By \ref{6.1.15} and \ref{prop:6.1.16},
    any non-zero $F \in K[T]$ can be decomposed as 
    $$F(T) = c(F) P_1(T) \cdots P_n(T),$$
    where $P_1,\cdots, P_n \in M$.
    It is unique up to a permutation.
    For any $P \in M$, denote by $\ord_P(F)$ the number of occurrence of $P$ in the decomposition of $F$.
    $$F(T) = c(F) \prod_{p \in M} P(T)^{\ord_P(F)}.$$
    If $F,G \in K[T]$ to be non-zero,
    then for any $P \in M$,
    $$\ord(F \cdot G) = \ord_P(F) + \ord_P(G).$$
    $$\ord_P(F + G) \ge \min\{\ord_P(F), \ord_P(G)\}.$$
    By convention, $\ord_P(0) = +\infty$.
    Moreover, $F \mid G$ if and only if
    $$\forall P \in M,\ \ord_P(G) \ge \ord_P(F).$$
    
    \quad Conversely, assume that any $I$ in $R$ is finitely generated.
    Let
    $$I_0 \subseteq I_1 \subseteq \cdots I_n \subseteq I_{n+1} \subseteq \cdots$$
    be an increasing sequence of ideals.
    Then $\dis I \coloneq I_n$ is an ideal in $R$, so $I = (x_1,\cdots,x_p)$.
    Let $n \in \NN$ such that 
    $$I_n \supseteq \{x_1,\cdots x_p\},$$
    hence $(x_1,\cdots,x_p) = I_n = I_{n+1} = \cdots$.
    Then any ideal $I \neq R$ is contained in some maximal ideal of $R$.
\end{remark}

\begin{propositionenv}
    Let $A$ be a commutative unitary ring.
    The set of all ideals of $A$ equipped with the relation of inclusion forms an order complete totally ordered set.
\end{propositionenv}
\begin{proofenv}
    Let $\Theta$ be a family of all ideals of $A$.
    If $\Theta$ is not empty, then $\dis \bigcap_{I \in \Theta} I$ is an ideal in $\Theta$.
    It is the the infimum of $\Theta$.
    Moreover, the ideal generated by $\dis \bigcup_{I \in \Theta } I$ is rhe supremum of $\Theta$.
    We deduce it by
    $$\sum_{I \in \Theta} I.$$
\end{proofenv}

\begin{definitionenv}
    Let $K$ be a field, $n \in \NN$, and $F_1,\cdots, F_n \in K[T]$.
    We call the \textbf{greatest common divisor} of $F_1,\cdots, F_n$ the unique monic polynomial that generates the ideal
    $$\sum_{i=1}^n F_i(T)K[T]$$
    We denote it $\gcd(F_1,\cdots, F_n)$.
    There exists $G_1,\cdots G_n \in K[T]$ such that
    $$\gcd(F_1,\cdots, F_n)(T) = \sum_{i=1}^n F_i(T)G_i(T).$$
    We call the \textbf{least common multiple} the polynomial that generates
    $$\bigcap_{i=1}^n F_i(T)K[T]$$
    We denote it $\lcm(F_1,\cdots, F_n)$.
\end{definitionenv}

\begin{propositionenv}
    Let $K$ be a field, $n \in \NN_{>0}$, $F_1, \cdots F_n \in K[T]$.
    For any irreducible monic $P \in K[T]$,
    $$\ord_P(\gcd(F_1,\cdots, F_n)) = \min_{i \in \{1,\cdots,n\}} \ord_{P}(F_i),$$
    $$\ord_P(\lcm(F_1,\cdots, F_n)) = \max_{i \in \{1,\cdots,n\}} \ord_{P}(F_i).$$
\end{propositionenv}
\begin{proofenv}
    By definition, for any $i \in \{1,\cdots,n\}$,
    $$\gcd(F_1,\cdots, F_n) \mid F_i,\quad F_i \mid \lcm(F_1,\cdots, F_n).$$
    Hence,
    $$\ord_P\left(\gcd(F_1, \cdots F_n)\right) \le \ord_P(F_i) \quad \text{and} \quad \ord_P\left(\lcm(F_1, \cdots F_n)\right) \ge \ord_P(F_i).$$
    Let $M$ be the set of monic irreducible $P \in K[T]$.
    $$F = \prod_{P \in M} P^{\min\{\ord_P(F_1),\cdots,\ord_P(F_n)\}}.$$
    For any $i \in \{1,\cdots,n\}$, $F \mid F_i$, so $F \mid \gcd(F_1,\cdots,F_n)$,
    $$\min\{\ord_{P}(F_1),\cdots, \ord_{P}(F_n)\} \le \ord_{P}\left(\gcd(F_1,\cdots,F_n)\right).$$
    Let
    $$G = \prod_{P \in M} P^{\max\{\ord_P(F_1),\cdots,\ord_P(F_n)\}}.$$
    For any $i \in \{1,\cdots,n\}$, $G$ is a multiple of $F_i$,
    hence
    $$\lcm(F_1,\cdots,F_n) \mid G.$$
    Therefore,
    $$\max\{\ord_P(F_1),\cdots, \ord_P(F_n)\} \ge \ord_{P}\left(\lcm(F_1,\cdots,F_n)\right).$$
\end{proofenv}


\section{Cayley-Hamilton Theorem}


\begin{definitionenv}\label{def:6.2.1}
    Let $K$ be a commutative unitary ring.
    A $K$-algebra is a $K$-module $A$ with a composition law
    $$\begin{array}{rcl}
        A \times A & \to & A\\
        (a,b) & \mapsto & ab
    \end{array}$$
    such that $A$ with the addition and the composition law forms a unitary ring such that
    $$\forall(\lambda,a,b) \in K \times A \times A,\ \lambda(ab) = \lambda(a)b = a\lambda(b).$$
    Let $A$ be an $K$ algebra and $V$ be a $K$-module.
    We equip $A \otimes V$ with the following structure of left $A$-module
    $$\begin{array}{rcl}
        A \times (A \otimes V) & \longrightarrow & A \otimes V\\
        (a,(b \otimes x)) & \longmapsto & ab \otimes x.
    \end{array}$$
\end{definitionenv}

\begin{propositionenv}\label{prop:6.2.2}
    Let $K$ be a commutative unitary ring,
    $A$ be a $K$ algebra and $V$ be a free $K$-module of finite rank $n$.
    Let $(e_i)_{i=1}^{n}$ be a basis of $V$.
    Then $(1 \otimes e_i)_{i=1}^{n}$ is a basis of $A \otimes V$.
\end{propositionenv}
\begin{proofenv}
    Let $M$ be a left $A$-module.
    One has the following bijections, which are functorial with respect to $M$.
    $$\begin{aligned}
        \Hom_A (A \otimes V, M) \cong & \{f : A \times V \rightarrow M \mid f \text{ is $K$-bilinear, } f (a,x) = a f(1,x)\}\\
        \cong & \Hom_K (V,M) \cong M^{\{e_1,\cdots,e_n\}} \cong M^{\{1 \otimes e_1,\cdots,1 \otimes e_n\}}.
    \end{aligned}$$
    Therefore, $A \otimes V$ is a free $A$-module generated by $\{1 \otimes e_i\}_{i=1}^{n}$.
\end{proofenv}

\begin{definitionenv}\label{def:6.2.3}
    Let $K$ be a commutative unitary ring,
    $V$ be a $K$-module and $\varphi \in \End_{K}(V)$.
    For any polynomial
    $$F(T) = \sum_{i=0}^{d} a_i T^i \in K[T],$$
    we denote by $F(\varphi) \in \End_{K}(V)$ the following $K$-linear endomorphism
    $$F(\varphi) = \sum_{i=0}^{d} a_i \varphi^i,\quad \varphi^0 = \Id_V.$$
    Note that $F \rightarrow F(\varphi)$ defines a homomorphism of $K$-algebras from $K[T]$ to $\End_{K}(V)$.
    Since $K[T]$ is commutative, we obtain that, for any $FG \in K[T]$, 
    $$F(\varphi) \circ G(\varphi) = G(\varphi) \circ F(\varphi).$$
    We equip $V$ with the structure of left $K[T]$-module in the following way
    $$\begin{array}{rcl}
        K[T] \times V & \longrightarrow & V\\
        (F(T),x) & \longmapsto & F(\varphi)(x).
    \end{array}$$
    This mapping is $K$-linear and hence defines a $K$-linear mapping $K[T] \otimes V \rightarrow V$ that sends $F \otimes x$ to $F(\varphi)(x)$.
    Note that if we consider $K[T] \otimes V$ as a $K[T]$-module and consider the above structure of $K[T]$-module on $V$,
    then the $K$-linear mapping $K[T]\otimes V \ni F \otimes x \mapsto F(\varphi)(x) \in V$ is actually $K[T]$-linear.
    Indeed, for any $F,G \in K[T]$, one has
    $$(FG)(\varphi)(x) = (F(\varphi) \circ G(\varphi))(x) = F(\varphi) (G(\varphi))(x).$$
\end{definitionenv}

\begin{lemmaenv}\label{lem:6.2.4}
    Let $K$ be a commutative unitary ring,
    $V$ be a free $K$-module of finite rank $n$ and $\varphi \in \End_{K}(V)$.
    Then one has
    $$(\varphi^a)^\vee = (\varphi^\vee)^a.$$
\end{lemmaenv}
\begin{proofenv}
    Let $(e_i)_{i=1}^{n}$ be a basis of $V$,
    $(e_i^\vee)_{i=1}^{n}$ be its dual basis.
    One has 
    $$e_1^\vee \wedge \varphi^\vee (e_2^\vee) \wedge \cdots \wedge \varphi^\vee (e_n^\vee) = e_1^\vee \wedge (e_2^\vee \circ \varphi) \wedge \cdots \wedge (e_n^\vee \circ \varphi).$$
    Therefore,
    $$\begin{aligned}
        &(e_1^\vee \wedge \varphi^\vee (e_2^\vee) \wedge \cdots \wedge \varphi^\vee (e_n^\vee)) (e_1\wedge \cdots \wedge e_n)\\
        = & \sum_{\sigma \in \mathfrak{S}_n} \sgn(\sigma) e_1^\vee(e_{\sigma(1)})\wedge e_2^\vee(\varphi(e_{\sigma(2)})) \wedge \cdots \wedge e_n^\vee(\varphi(e_{\sigma(n)})).
    \end{aligned}$$
    We obtain
    $$\begin{aligned}
        & (e_1^\vee \wedge \varphi^\vee(e_2^\vee) \wedge \cdots \wedge \varphi^\vee(e_n^\vee))(e_1 \wedge \cdots \wedge e_n)\\
        = & (e_1^\vee \wedge \cdots \wedge e_n^\vee) (e_1 \wedge \varphi(e_2) \wedge \cdots \wedge \varphi(e_n))\\
        = & (e_1^\vee \wedge \cdots \wedge e_n^\vee)(\varphi^a(e_1) \wedge e_2 \wedge \cdots \wedge e_n)\\
        = & \sum_{\sigma \in \mathfrak{S}_n} \sgn(\sigma) e^\vee_{\sigma(1)} (\varphi^a(e_1))e_{\sigma(2)^\vee}(e_2) \cdots e_{\sigma(n)}^\vee(e_n)\\
        = & (e_1^\vee \circ \varphi^a) \wedge e_2^\vee \wedge \cdots \wedge e_n^\vee (e_1 \wedge \cdots \wedge e_n).
    \end{aligned}$$
    We obtain
    $$(\varphi^a)^\vee (e_1^\vee) = (\varphi^\vee)^a (e_1^\vee).$$
\end{proofenv}

\begin{lemmaenv}\label{lem:6.2.5}
    Let $K$ be a commutative unitary ring,
    $V$ be a free $K$-module of finite rank $n$ and $\varphi, \psi \in \End_{K}(V)$.
    If $\varphi^\vee = \psi^\vee$, then $\varphi = \psi$.
\end{lemmaenv}
\begin{proofenv}
    Let $(e_i)_{i=1}^{n}$ be a basis of $V$,
    $(e_i^\vee)_{i=1}^{n}$ be its dual basis.
    By assumption $\varphi^\vee = \psi^\vee$, we have
    $$ e_i^\vee \circ \varphi = e_i^\vee \circ \psi,\; i \in \{1,\cdots,n\}.$$
    For any $x \in V$,
    $$x = \sum_{i=1}^{n} e_i^\vee(x) \cdot e_i.$$
    Therefore, for any $y \in V$,
    $$\varphi(y) = \sum_{i=1}^{n} (e_i^\vee \circ \varphi) (y) = \sum_{i = 1}^{n} (e_i^\vee \circ \psi) (y) = \psi(y).$$
\end{proofenv}

\begin{propositionenv}\label{prop:6.2.6}
    Let $K$ be a commutative unitary ring,
    $V$ be a free $K$-module of finite rank $n$ and $\varphi \in \End_{K}(V)$.
    Then
    $$\varphi \circ \varphi^a = \det(\varphi) \Id_V.$$
\end{propositionenv}
\begin{proofenv}
    Applying the theorem \ref{thm:4.7.1} to $\varphi^\vee$, by lemma \ref{lem:6.2.4} and proposition \ref{thm:LinearMap.det_dualMap}, 
    $$(\varphi \circ \varphi^a)^\vee = (\varphi^a)^\vee \circ \varphi^\vee = \det(\varphi^\vee) \Id_{V^\vee} = \det(\varphi) \Id_{V^\vee}.$$
    Hence, by lemma \ref{lem:6.2.5}, we obtain the claimed equality.
\end{proofenv}

\begin{theoremenv}[Cayley-Hamilton]
    Let $K$ be a commutative unitary ring,
    $V$ be a free $K$-module of finite rank $n$ and $\varphi \in \End_{K}(V)$.
    Let $T \cdot \Id_{V} - \varphi$ denote the $K[T]$-linear endomorphism
    $$T \otimes \Id_{V} - 1 \otimes \varphi$$
    of $K[T] \otimes V$ that sends $F(T) \otimes x$ to 
    $$TF(T) \otimes x - F(T) \otimes \varphi(x),$$
    and let
    $$P_\varphi(T) = \det(T \cdot \Id_{V} - \varphi).$$
    Then
    $$P_\varphi(\varphi) = 0$$
    in $\End_{K}(V)$.
\end{theoremenv}
\begin{proofenv}
    For any element of the form $F(T) \otimes x \in K[T] \otimes V$,
    the image of 
    $$(T\cdot \Id_{V} - \varphi)(F(T) \otimes x)$$
    in $V$ by the $K[T]$-linear mapping $K[T] \otimes V \rightarrow V$
    described in definition \ref{def:6.2.3} vanishes, cause
    $$\varphi \circ F(\varphi) = F(\varphi) \circ \varphi.$$
    Indeed,
    $$(T\Id_{V} - \varphi)(F \otimes x) = T (F \otimes x) - F \otimes \varphi(x),$$
    $$\mu(T \Id_{V} - \varphi)(F \otimes x) = \varphi F (\varphi)(x) - F(\varphi)(\varphi(x)).$$
    By linearity, we deduce that for any $\xi \in K[T] \otimes V$,
    the image of $(T \Id_{V} - \varphi) (\xi)$ by the $K[T]$-linear mapping
    $$\begin{array}{rrcl}
        \mu : & K[T] \otimes V & \longrightarrow & V \\
        & F(T) \otimes x & \longmapsto & F(\varphi)(x)
    \end{array}$$
    vanishes.
    By proposition \ref{prop:6.2.6}, 
    $$(T - \varphi) \circ (T - \varphi)^a = P_{\varphi}(T) \Id_{K[T] \otimes V}.$$
    Let $x \in V$, consider $1 \otimes x$, then
    $$P_{\varphi} (T) \otimes x = P_{\varphi}(T)(1 \otimes x) = (T\Id_{V} - \varphi)(\xi),$$
    where $\xi = (T\Id_{V} - \varphi)^a(1 \otimes x)$.
    Taking the image of this under $\mu$, we get
    $$P_{\varphi}(\varphi)(x) = 0.$$
\end{proofenv}

\begin{definitionenv}
    Let $K$ be a commutative unitary ring,
    $V$ be a free $K$-module of finite rank $n$ and $\varphi \in \End_{K}(V)$.
    The polynomial $P_\varphi(T) = \det(T \cdot \Id_{V} - \varphi)$ is called the \textbf{characteristic polynomial} of $\varphi$.

    \quad Let $(e_i)_{i=1}^n$ be a basis of $V$ over $K$.
    Then $(1 \otimes e_1)_{i=1}^{n}$ forms a basis of the $K[T]$-module $K[T] \otimes V$.
    By abuse of notation, for any $x \in V$,
    we still use the notation $x$ to denote the element $1 \otimes x$ in $K[T] \otimes V$.
    We therefore identify $V$ as a sub-$K$-module of $K[T] \otimes V$ with this notation
    $$(T\Id_{V} - \varphi)(e_i) = Te_i - \varphi(e_i).$$
    Therefore,
    $$(T\Id_{V} - \varphi)(e_1) \wedge \cdots \wedge (T\Id_{V} - \varphi)(e_n) = (Te_1 - \varphi(e_1)) \wedge \cdots \wedge (Te_n - \varphi(e_n)).$$
    This can be written as 
    $$\begin{aligned}
        &T^n(e_1 \wedge \cdots \wedge e_n)\\
        & - T^{n-1} \sum_{i=1}^{n}(e_1 \wedge \cdots \wedge e_{i-1} \wedge \varphi(e_i) \wedge e_{i+1} \wedge \cdots \wedge e_n) + \cdots \\
        &+ (-1)^n \varphi(e_1) \wedge \cdots \wedge \varphi(e_n).
    \end{aligned}$$
    We see that $P_{\varphi}$ is a monic polynomial of degree $n$,
    and the coefficient of the polynomial is $(-1)^n \det(\varphi)$.
    We denote by $\tr(\varphi)$ the opposite of the coefficient of $T^{n-1}$ in $P_{\varphi}$,
    namely,
    $$(\tr (\varphi))(e_1 \wedge \cdots \wedge e_n) = \sum_{i=1}^{n} (e_1 \wedge \cdots \wedge e_{i-1} \wedge \varphi(e_i) \wedge e_{i+1} \wedge \cdots \wedge e_n).$$
    We can write $P_{\varphi} (T)$ as
    $$P_{\varphi}(T) = T^n - s_1 (\varphi) T^{n-1} + \cdots + (-1)^{n} s_n (\varphi).$$ 
    In particular, $\tr(\varphi) = s_1 (\varphi)$ and $\det(\varphi) = s_n(\varphi).$
\end{definitionenv}

\begin{remark}
    Assume that
    $$P_{\varphi} (T) = (T - \lambda_1) \cdots (T - \lambda_n),$$
    then
    $$s_{l} = \sum_{1 \le i_1 \le \cdots \le i_{l} \le n} \lambda_{i_1} \cdots \lambda_{i_{l}}.$$
    By definition, for any $(x_{1},\cdots,x_{n}) \in V^n$,
    $$(Tx_1 + \varphi(x_1)) \wedge \cdots \wedge (Tx_n + \varphi(x_n)) = (T^n + s_1(\varphi) T^{n-1} + \cdots + s_d(\varphi)) x_1 \wedge \cdots \wedge x_{n}.$$
    Hence,
    $$s_l (\varphi)(x_1 \wedge \cdots \wedge x_n) = \sum_{\substack{I = (i_1,\cdots,i_{l})\\1\le i_1 \le \cdots \le i_{l} \le n}} z_1^I \wedge \cdots \wedge z_n^I,$$
    where 
    $$z_j^I = 
    \begin{cases}
        x_j, \ &j \notin I\\
        \varphi(x_j), \ &j \in I
    \end{cases}.$$
    Let $F_{l} : \End_{K}(V)^l \rightarrow K$ be the multilinear mapping
    $$F_{l}(\varphi_1, \cdots, \varphi_l) x_1 \wedge \cdots x_d = \sum_{\substack{I = (i_1,\cdots,i_{l})\\1\le i_1 \le \cdots \le i_{l} \le n}} u_1^I \wedge \cdots \wedge u_d^I,$$
    where 
    $$u_j^I = 
    \begin{cases}
        x_j \ &j \notin I\\
        \varphi_k(x_j) \ &j \in I
    \end{cases}.$$
    $F_l$ is multilinear form and $s_l (\varphi) = F_{l}(\varphi,\cdots,\varphi)$.
    Hence $s_l$ is a homogenous polynomial of degree $l$ on $\End_{K}(V)$.
\end{remark}


\section{Reduction of Linear Endomorphisms}


In this section, we fix a field $K$.

\begin{definitionenv}\label{def:6.3.1}
    Let $V$ be a finite dimensional vector space over $K$.
    If $\varphi$ is a $K$-linear endomorphism of $V$,
    we let $\lambda_{\varphi}$ to be the homomorphism of $K$-algebras
    $$\begin{array}{rrcl}
        \lambda_{\varphi} : & K[T] & \longrightarrow & \End_{K}(V)\\
        & F(T) & \longmapsto & F(\varphi).
    \end{array}$$
    By Cayley-Hamilton theorem, $P_{\varphi} \in \ker(\lambda_{\varphi})$.
    As a consequence, $\ker (\lambda_{\varphi})$ is a non-zero ideal in $K[T]$ (a principal ideal domain).
    We denote by $Q_{\varphi}(T)$ the unique monic polynomial that generates $\ker(\lambda_{\varphi})$.
    
    \quad Note that $F \in K[T]$ satisfies $F(\varphi) = 0$ if and only if $Q_{\varphi} \mid F$.
    We call $Q_\varphi$ the \textbf{minimal polynomial} of $\varphi$.
\end{definitionenv}

\begin{propositionenv}\label{prop:6.3.2}
    Let $V$ be a finite dimensional vector space over $K$,
    $\varphi \in \End_{K}(V)$,
    and $W$ be a vector subspace of $V$ such that $\varphi(W) \subseteq W$.
    Let $\psi : W \rightarrow W$ be the restriction of $V$ to $W$.
    Then $Q_{\psi} \mid Q_{\varphi}$.
\end{propositionenv}
\begin{proofenv}
    $Q_{\varphi}(\varphi) = 0$, hence, for any $x \in W$,
    $$Q_{\varphi}(\psi)(x) = Q_{\varphi}(\varphi)(x) = 0.$$
    Hence, $Q_{\varphi}(\psi) = 0$, then $Q_{\psi} \mid Q_{\varphi}$.
\end{proofenv}

\begin{definitionenv}\label{def:6.3.3}
    Let $V$ be a vector space over $K$,
    $I$ be a set and $(V_i)_{i \in I}$ be a family of vector subspaces of $V$.
    Recall that, in the case where the $K$-linear mapping
    $$ \bigoplus_{i \in I} V_i \longrightarrow V, \quad (x_i)_{i \in I} \longmapsto \sum_{i \in I} x_i$$
    is an isomorphism of $K$-vector spaces, we say that $V$ is a direct sum of $(V_i)_{i \in I}$.
    
    \quad Note that if for any $i \in I$, $(e_{i,j})_{j \in J_i}$ is a basis of $V_i$, then
    $$e_{i,j}, \quad i \in I, j \in J_i$$
    forms a basis of $V$ over $K$.
\end{definitionenv}

\begin{propositionenv}\label{prop:6.3.4}
    Let $V$ be a finite dimensional vector space over $K$,
    $\varphi \in \End_{K}(V)$,
    $V_1,\cdots,V_n$ be vector subspaces of $V$.
    We assume that, for any $i \in \{1,\cdots,n\}$, $\varphi(V_i) \subseteq V_{i}$.
    Let $\varphi_i:V_i \rightarrow V_i$ be the restriction of $\varphi$ to $V_i$.
    \begin{enumerate}
        \item If $V = V_1 + \cdots + V_n$, then $Q_{\varphi} (T) = \lcm(Q_{\varphi_1}(T), \cdots, Q_{\varphi_n}(T))$.
        \item If $V$ is the direct sum of $V_1,\cdots,V_n$, then 
            $$P_{\varphi}(T) = P_{\varphi_1}(T) \cdots P_{\varphi_n}(T).$$
    \end{enumerate}
\end{propositionenv}
\begin{proofenv}
    \quad
    \begin{enumerate}
        \item Let $G = \lcm(Q_{\varphi_1}, \cdots, Q_{\varphi_n})$.
            For any $i \in \{1,\cdots,n\}$,
            let $F_i \in K[T]$ be such that $G = F_{i} Q_{\varphi_i}$.
            Let $x \in V$, $x = x_1 + \cdots + x_n$, $x_i \in V_i$, $i \in \{1,\cdots,n\}$.
            One has
            $$G(\varphi)(x) = \sum_{i=1}^{n} G(\varphi)(x_i) = \sum_{i=1}^{n} G(\varphi_i)(x_i) = \sum_{i=1}^{n} F_i (\varphi_i) Q_{\varphi_i}(\varphi_i)(x).$$
            Hence, $G$ is a vanishing polynomial of $\varphi$, so $Q_\varphi \mid G$.
            Conversely, by proposition \ref{prop:6.3.2}, one ahs that $Q_{\varphi_i} \mid Q_{\varphi}$ for any $i \in \{1,\cdots,n\}$.
            Hence, $G \mid Q_{\varphi}$.
        \item For any $i \in \{1,\cdots,n\}$, let $(e_{i,j})_{j=1}^{d_i}$ be a basis of $V_i$, $d_i = \dim(V_i)$.
            One has
            $$(Te_{i,1} - \varphi(e_{i,1})) \wedge \cdots \wedge (Te_{i,d_i} - \varphi(e_{i,d_i})) = P_{\varphi_i} (T) e_{i,1} \wedge \cdots \wedge e_{i,d_i}.$$
            $$P_{\varphi}(T) = P_{\varphi_{1}}(T) \cdots P_{\varphi_{n}}(T).$$
    \end{enumerate}
\end{proofenv}

\begin{theoremenv}
    Let $V$ be a finite dimensional vector space over $K$,
    let $F_1,\cdots,F_n$ be non-zero polynomials in $K[T]$ and $F = F_1 \cdots F_n$.
    Assume that $F_1,\cdots F_n$ are pairwise coprime,
    that is for any $i \neq j \in \{1,\cdots,n\}$, $F_i,F_j$ do not have any non-trivial common irreducible factor.
    Then $\ker(F(\varphi))$ is the direct sum of
    $$\ker(F_1(\varphi)), \cdots , \ker(F_n(\varphi)).$$
\end{theoremenv}
